\section{Grundlagen}
\label{sec:grundlagen}

Tcpdump\footnote{\url{http://www.tcpdump.org/}} ist eine Software um
Netzwerkpakete aufzuzeichnen. Tcpdump benutzt die Bibliothek libpcap, welche
eine einheitliche API anbieten um Netzwerkpakete auf unterschiedlichen
Betriebssystemen aus dem Netzwerkverkehr herauszufiltern. Ziel hierbei ist es
Paket möglichst früh im Netzwerkstack zu filtern, wodurch nur Pakete
von Interesse weiterverarbeitet werden müssen.

Diese Kapitel beschreibt wie tcpdump/libpcap in Linux Netzwerkpakete filtert.
Dazu wird wie in
\cite{Insolvibile:2001:KKL,Insolvibile:2002:KKIa,Insolvibile:2002:KKIb}
beschrieben der Verlauf eines eingehenden TCP Paketes durch den Linux Kernel
betrachtet.


\begin{figure} \centering
\includegraphics[width=.7\textwidth]{images/network-stack.eps}
\caption{Schematischer Paketfluss beim Empfang eines TCP Paketes im Linux
Kernel}\label{fig:network-stack}
\end{figure}

In Abbildung~\ref{fig:network-stack} sind die wichtigsten Teile des Linux
Networkstacks dargestellt, welche bei dem Empfang eines TCP Paketes eine Rolle
spielen. Im Userspace wird ein Webserver (nginx) betrieben, welcher über einen
TCP Socket HTTP Requests empfängt. Parallel läuft tcpdump/libpcap welches einen
speziellen \texttt{PF\_PACKET} Socket nutzt um alle empfangen Pakete zu
untersuchen (siehe Listing~\ref{lst:pfpacket}).

\begin{lstlisting}[caption={Erzeugen eines \texttt{PF\_PACKET} Sockets},label=lst:pfpacket]
socket(PF_PACKET, SOCK_RAW, htons(ETH_P_ALL));
\end{lstlisting}


Ein solcher Socket ermöglicht es Pakete direkt von der Netzwerkkarte zu
erhalten ohne das die Pakete den normalen Netzwerkstack durchlaufen, dies ist
wichtig für tcpdump da es das komplette Paket aufzeichnen will.

Erreicht ein TCP Paket die Netzwerkkarte wird es zuerst in eine Queue in der
Netzwerkkarte (1) abgespeichert. Anschließend wird das Paket in die zuständige
CPU Eingangsqueue übergeben (2). Danach übernimmt die Kernel Funktion
\texttt{net\_rx\_action} die Weitergabe des Paketes. Dazu werden zwei Listen
von Handler betrachtet. Als erstes wird das Paket an alle Handler übergeben,
welche alle Pakete empfangen. Dazu zählt auch der \texttt{PF\_PACKET} Socket
von tcpdump/libpcap. Anschließend werden spezielle Handler für das Paket
aufgerufen. In dem Fall eines TCP Paketes ist dies der IP Handler, welcher den
IP Header liest, auswertet und entfernt.  Anschließend wird das Paket an den
TCP Handler übergeben. Dieser wertet den TCP Header aus und reicht das Paket an
den zuständigen Socket weiter. Wenn ein Paket einen Socket erreicht wird es in
dessen Eingangsbuffer (3) abgelegt.

Wie in Abbildung~\ref{fig:network-stack} zu sehen ist, gibt es mindestens drei
Paketbuffer die ein Paket durchläuft. Wenn einer dieser Buffer voll ist, weil
bereits zu viele Paket angenommen aber nicht verarbeitet wurden, wir ein neues
Paket verworfen. Unter der Fragestellung unter welchen Umständen tcpdump Pakete
verwirft, welcher allerdings die Anwendung (nginx) erreichen, sind die ersten
beiden Buffer (NIC und CPU Queue) nicht interessant. Ein Paket, welches bereits
an dieser Stelle verworfen wird, erreicht weder die Anwendung noch tcpdump.
Insofern ist nur der dritte Buffer vom Socket interessant. An dieser Stelle hat
das Paket bereits zwei unterschiedliche Wege durchlaufen. Dabei muss die Anzahl
der Paket in beiden Sockets nicht identisch sein, da der \texttt{PF\_PACKET}
generell alle Pakete erhält. Dies bedeutet, wenn neben dem HTTP Verkehr noch
weiterer Netzwerkverkehr existiert wird dieser ebenfalls an den tcpdump/libpcap
Socket weitergeleitet.

Daher ist es nötig diesen Traffic möglichst frühzeitig zu filtern. Eine
Filterung im Userspace würde das Problem nicht beheben, da immer noch alle
Pakete erst einmal den Socket erreichen. Linux stellt für diese Zweck den Linux
Socket
Filter\footnote{\url{https://www.kernel.org/doc/Documentation/networking/filter.txt}}
(LSF) bereit. Diese API wurde mit Kernel 2.2 einführt und ermöglicht es ein Filter
an einem Socket zu registrieren (siehe Listing~\ref{lst:attach}).

\begin{lstlisting}[caption={Anhängen eines Filters an eine Socket},label=lst:attach]
setsockopt(socket, SOL_SOCKET, SO_ATTACH_FILTER, filter, sizeof(*filter));
\end{lstlisting}

Der Filter ist ein Programm in der Assembler ähnlichen Sprache Berkley Paket
Filter (BPF)~\cite{DBLP:conf/usenix/McCanneJ93}. Die Sprache umfasst einfache
Vergleich- und Rechenoperation und arbeitet direkt auf den Daten des
Netzwerkpaketes. Dadurch kann der Filter direkt auf einzelne Header Felder
der verschieden Protokolle zugreifen.\ tcpdump/libpcap erlaubt es logische
Ausdrücke zum erstellen eines Filters in BPF umzuwandeln. Das Beispiel
in Listing~\ref{lst:filter} zeigt den BPF Code für einen Filter der alle
TCP Pakete für die Ziel-IP \texttt{10.3.9.21} und Ziel-Port \texttt{80} akzeptiert.

\begin{lstlisting}[numbers=none,caption={BPF for \texttt{tcpdump -d ip dst host 10.3.9.21 and tcp dst port 80}},label={lst:filter}]
(000) ldh      [12]
(001) jeq      #0x800           jt 2    jf 12
(002) ld       [30]
(003) jeq      #0xa030915       jt 4    jf 12
(004) ldb      [23]
(005) jeq      #0x6             jt 6    jf 12
(006) ldh      [20]
(007) jset     #0x1fff          jt 12   jf 8
(008) ldxb     4*([14]&0xf)
(009) ldh      [x + 16]
(010) jeq      #0x50            jt 11   jf 12
(011) ret      #65535
(012) ret      #0
\end{lstlisting}

Das BPF Programm lädt Daten aus dem Netzwerkpaket und vergleicht diese mit dem
festgelegten Filter. Zuerst wird der EtherType aus dem Ethernet Header geladen
(000) und verglichen ob es sich um ein IPv4 Paket (Type: 0x0800) handelt (001).
Ist der Vergleich erfolgreich wird mit der nächsten Instruktion fortgefahren,
ansonsten wird zu letzten Instruktion (012) gesprungen. Als zweites wird die
Ziel-Adresse aus dem IP Header geladen (003) und mit der IP 10.3.9.21
(0xa030915) verglichen (004). Danach wird das Protokolfeld des IP Headers
geladen (004) und überprüft, dass es sich um ein TCP Paket (0x6) handelt (005).
In der folgende Instruktion wird sichergestellt, dass das IP Paket nicht
fragmentiert wurde oder es sich um das erste Fragment handelt, also der IP
Fragment Offset 0 ist (007). Anschließend wird die IP Headerlänge im
Hilfsregister X gespeichert (008), dies ist nötig, da der IP Header durch
Optionen eine variable Länge besitzt. Ein Zugriff auf den TCP Header ist nur
mit dieser Information möglich. Als letztes wird der Ziel-Port aus dem TCP
Header geladen (009) und mit Port 80 (0x50) verglichen (010). Wenn alle diese
Tests erfolgreich waren wird mit Instruktion (011) das Paket akzeptiert.
Ansonsten wird das Paket mit Instruktion (012) abgelehnt. Das BPF Programm gibt
die Anzahl der Bytes zurück, welche von dem Paket weitergeleitet werden sollen.
In libpcap ist der Standard dafür 65535 Bytes (siehe Instruktion 011), welches
die maximale Größe eines IP Paketes entspricht. Der Rückgabewert 0 gibt an
das kein Byte des Paketes weitergegeben werden soll (012).

Ist ein Filter für einen Socket registriert, wird dieser vom Handler durch die Kernel
Funktion \texttt{sk\_run\_filter} ausgeführt, bevor
das Paket an den Socket übergeben wird. Dadurch kann die Anzahl der Paket,
welche den tcpdump/libpcap Socket erreichen effizient reduziert werden.

Der tcpdump/libpcap Socket Buffer wird als Memory Mapped Packet Ring
Buffer\footnote{\url{https://www.kernel.org/doc/Documentation/networking/packet_mmap.txt}}
angelegt. Dies ermöglicht es effizienter auf Paket im Userspace zu zugreifen.
Dies verhindert, dass jedes Paket einzeln durch einen \texttt{recv} Aufruf
empfangen werden muss.  Die Größe des Buffers ist standardmäßig 2~MB\@. Der
Buffer wird in Frames aufgeteilt, welche es ermöglichen das aufgezeichnete
Paket und Metadaten zu speichern. Dabei richtet sich die Framegröße nach der
Anzahl der Bytes die pro Paket aufgezeichnet werden.

Es zeigt sich, dass es mindestens drei Merkmale gibt, welche beeinflussen wieviel
Paket tcpdump/libpcap aufzeichnen kann. Dies ist zu erst der Filter welcher
möglichst viele Paket bereits im Kernel verwerfen sollte. Anschließend die
Anzahl der Bytes die pro gefiltertem Paket weitergegeben werden und die
Buffergröße die verwendet wird um die Paket vorzuhalten. Die Auswirkung
dieser Aspekte soll in dieser Arbeit mit Hilfe von Messungen untersucht werden.
Außerdem ist bereits aus Abbildung~\ref{fig:network-stack} ersichtlich das die
Ausführung von tcpdump eine Auswirkung auf andere Services hat, da jedes Paket
nun von einem zusätzlichen Handler betrachtet und gefiltert wird. Daher hat die
Komplexität und Effektivität des tcpdump Filters einen direkten Einfluss auf
die Geschwindigkeit mit der Paket verarbeitet werden.
